\begin{table}[h!]
    \caption[Semantic Field of Each Function-Attitude. \verified]{Semantic Field of Each Function-Attitude.\\
    \verified~\textcolor{gray}{\small Source: Page 4-7, Section ``The Eight Function-Attitudes Unpacked'' from Beebe~\cite{beebe}.}}
    \begin{center}

        \begin{tabular}{ccccc}
            \toprule
            \multirow{2}{*}{\textbf{Function-attitude}} & \multirow{2}{*}{\textbf{Full Name}} & \multicolumn{3}{c}{\textbf{Semantic Field (Triangle)}} \tabularnewline
            \cmidrule(lr){3-5}
            & & \small{人格面具 (Persona)} & \small{自我 (Ego)} & \small{自性 (Self)} \tabularnewline
            \midrule

            \rowcolor{black!10}Se & Extraverted Sensing & Engaging & Experiencing & Enjoying \tabularnewline
            Si & Introverted Sensing & Implementing & Verifying & Accounting \tabularnewline

            \rowcolor{black!10}Ne & Extraverted Intuition & Entertaining & Envisioning & Enabling \tabularnewline
            Ni & Introverted Intuition & Imagining & Knowing & Divining \tabularnewline

            \rowcolor{black!10}Te & Extraverted Thinking & Regulating & Planning & Enforcing \tabularnewline
            Ti & Introverted Thinking & Naming & Defining & Understanding \tabularnewline

            \rowcolor{black!10}Fe & Extraverted Feeling & Validating & Affirming & Relating \tabularnewline
            Fi & Introverted Feeling & Judging & Appraising & Establishing the Value \tabularnewline

            \bottomrule
        \end{tabular}
    \end{center}

    {%\small
    \begin{description}

        \item[Semantic Field]
            The middle word is the apex of a triangle.
        The first of the three words (as we read across for each of the eight mental processes)
        can be thought of as what the process looks like on the surface, to another person seeing someone for the first time using that process -- what we might call its \textit{persona} level.
        The second, central, word captures the heart of the process -- the one that is
        embraced by the \textit{ego} as its chief concern.
        The third word embodies what the process becomes at its most evolved level when it
        is working in sympathetic harmony with what Jung call the \textit{Self}:
        one might call this the goal of the process.

        \item[Example] Introverted feeling (Fi) can develop from an initial, near instinctual ‘judging’ into a mature, reflective ‘appraising’ before evolving to a more far-seeing ‘establishing the value.’ In this latter stage, which aspires to wisdom, the feeling reaction finds its ground in a more ideal, archetypal realm, which allows it to discriminate the value that has led to the earlier judgments and appraisals.

        \item[Suggestion to Learners]
            To focus on the central words in my horizontal sequences for the different functions of consciousness, using the wing words as the beginning and the end of the process whose functioning they are trying to become conscious of.
            Using this method of observation, it will become clearer the degree to which some people are centrally engaged in appraising, defining, knowing, or verifying as \ul{the internal process by which they hope to master the objects they encounter} (or say, internal methods by which they have already learned to deal with what they encounter).
            Other people are more occupied with affirming, planning, envisioning, or experiencing as \ul{their way of expressing their willingness to subject themselves to the influence of the objects with which they are trying to connect} (or say, methods by which they enact a genuine encounter with the Other).
            Seeing this fundamental difference between the introverted and extraverted functions, we can gradually come to understand not just the ways the functions differ in their introverted and extraverted aspects, but what the different purposes are that make them functions.

    \end{description}
    }



\end{table}
